
\SGARefTarget{sga3:xx:expose}


\section{Elementary systems. The groups \(U_\alpha\) and \(U_{-\alpha}\)}
\SGARefTarget{sga3:xx:section:1}

\paragraph{\textbf{Recollection \SGARefLink{sga3:xx:section:1}{1.1}.}}
\SGARefTarget{sga3:xx:1.1}
\emph{Let \(S=\Spec(k)\), where \(k\) is an algebraically closed
field, and let \(G\) be a reductive \(S\)-group of semisimple rank
one and \(T\) a maximal torus of \(G\), necessarily split. Then}
\[
  \mathfrak g=\mathfrak t\oplus\mathfrak g^\alpha
  \oplus\mathfrak g^{-\alpha},
\]
\emph{where \(\alpha\) and \(-\alpha\) are the roots of \(G\) with
respect to \(T\). Moreover, there exist two group monomorphisms}
\[
  p_\alpha\colon\mathbb G_{a,S}\longrightarrow G,
  \qquad
  p_{-\alpha}\colon\mathbb G_{a,S}\longrightarrow G
\]
\emph{such that}
\[
  t\,p_\alpha(x)\,t^{-1}=p_\alpha\bigl(\alpha(t)x\bigr),
  \qquad
  t\,p_{-\alpha}(x)\,t^{-1}
    =p_{-\alpha}\bigl(\alpha(t)^{-1}x\bigr)
\]
\emph{for every \(S'\to S\), \(t\in T(S')\), and
\(x\in\mathbb G_a(S')\), and such that the morphism}
\[
  \mathbb G_{a,S}\times_S T\times_S\mathbb G_{a,S}
    \longrightarrow G,
  \qquad
  (y,t,x)\longmapsto p_{-\alpha}(y)\,t\,p_\alpha(x),
\]
\emph{is radicial and dominant \textup{(\emph{Bible}, Section 13.4,
Corollary 2 to Theorem 3)}.}

Since its tangent map at the identity is bijective, this morphism is
also separable, hence birational. Zariski's Main Theorem
\textup{(EGA III\(_1\), 4.4.9)} therefore shows that it is an open
immersion.

\paragraph{\textbf{Lemma 1.2.}}
\SGARefTarget{sga3:xx:1.2}
\emph{Let \(S\) be a scheme, \(G\) an \(S\)-group scheme, \(T\) a
torus of \(G\), \(Q\) a subtorus of \(T\), and \(\alpha\) a
character of \(T\) whose restriction to \(Q\) is nontrivial on every
fiber. Let
\[
  p_\alpha\colon\mathbb G_{a,S}\longrightarrow G
  \quad\textup{and}\quad
  p_{-\alpha}\colon\mathbb G_{a,S}\longrightarrow G
\]
be group morphisms normalized by \(T\), with multipliers \(\alpha\)
and \(-\alpha\), respectively. Suppose that the morphism
\[
  u\colon\mathbb G_{a,S}\times_S T\times_S\mathbb G_{a,S}
    \longrightarrow G,
  \qquad
  u(y,t,x)=p_{-\alpha}(y)\,t\,p_\alpha(x),
\]
defined on points by the displayed formula, is an open immersion.
Finally, let \(q\geq 0\) be an integer and let
\[
  p\colon\mathbb G_{a,S}\longrightarrow G
\]
be a group morphism such that, for every \(S'\to S\),
\(t\in Q(S')\), and \(x\in\mathbb G_a(S')\),
\[
  \intaut(t)^q\bigl(p(x)\bigr)=p\bigl(\alpha(t)x\bigr).
\]
Then there is a unique \(\nu\in\mathbb G_a(S)\) such that
\[
  p(x)=p_\alpha(\nu x^q).
\]}

Indeed, let \(\Omega\) be the image of \(u\), and put
\(U=p^{-1}(\Omega)\). This is an open subset of
\(\mathbb G_{a,S}\) containing the zero section. For every section
\(t\) of \(Q\), multiplication by \(\alpha(t)\) leaves \(U\)
globally invariant. We have \(U=\mathbb G_{a,S}\). It is enough to
check this when \(S\) is the spectrum of an algebraically closed
field \(k\); then
\[
  \alpha\colon Q(k)\longrightarrow k^\times
\]
is surjective, so \(U(k)\supset k^\times\), and hence
\(U=\mathbb G_{a,k}\).

There are consequently three morphisms
\[
  a\colon\mathbb G_{a,S}\longrightarrow\mathbb G_{a,S},
  \qquad
  b\colon\mathbb G_{a,S}\longrightarrow T,
  \qquad
  c\colon\mathbb G_{a,S}\longrightarrow\mathbb G_{a,S}
\]
such that
\[
  p(x)=p_{-\alpha}\bigl(a(x)\bigr)\,b(x)\,
       p_\alpha\bigl(c(x)\bigr).
\]
The condition on \(p\) becomes
\[
  a\bigl(\alpha(t)x\bigr)=\alpha(t)^{-q}a(x),\qquad
  b\bigl(\alpha(t)x\bigr)=b(x),\qquad
  c\bigl(\alpha(t)x\bigr)=\alpha(t)^q c(x).
\]
For the same reason, for every \(S'\to S\) and
\(z\in\mathbb G_m(S')\),
\[
  a(zx)=z^{-q}a(x),\qquad b(zx)=b(x),\qquad c(zx)=z^q c(x),
\]
and therefore
\[
  z^q a(z)=a(1),\qquad b(z)=b(1),\qquad c(z)=z^q c(1).
\]
Since \(\mathbb G_{m,S}\) is schematically dense in
\(\mathbb G_{a,S}\), for every \(S'\to S\) and
\(x\in\mathbb G_a(S')\) we obtain
\[
\begin{aligned}
  x^q a(x)&=a(1)=a(0)=0, &&\text{hence }a=0,\\
  c(x)&=x^q c(1)=\nu x^q,
    &&\text{for some }\nu\in\mathbb G_a(S),\\
  b(x)&=b(1)=b(0)=e, &&\text{hence }b=e.
\end{aligned}
\]
This completes the proof.

\paragraph{\textbf{Definition 1.3.}}
\SGARefTarget{sga3:xx:1.3}
\emph{An \(S\)-elementary system is a triple \((G,T,\alpha)\) in
which}
\begin{enumerate}
\item[\textup{(i)}] \emph{\(G\) is a reductive \(S\)-group of
  semisimple rank one \textup{(\SGARefLink{sga3:xix:definition:2:7}{Exposé XIX, 2.7})};}
\item[\textup{(ii)}] \emph{\(T\) is a maximal torus of \(G\);}
\item[\textup{(iii)}] \emph{\(\alpha\) is a root of \(G\) with
  respect to \(T\) \textup{(\SGARefLink{sga3:xix:definition:3-2}{Exposé XIX, 3.2})}.}
\end{enumerate}
Thus there is a direct-sum decomposition
\textup{(\SGARefLink{sga3:xix:lemma:3:5}{Exposé XIX, 3.5})}
\[
  \mathfrak g=\mathfrak t\oplus\mathfrak g^\alpha
  \oplus\mathfrak g^{-\alpha},
\]
where \(\mathfrak g^\alpha\) and \(\mathfrak g^{-\alpha}\) are
locally free of rank one.

\paragraph{\textbf{1.4.}}
\SGARefTarget{sga3:xx:1.4}
If \((G,T,\alpha)\) is an \(S\)-elementary system, then
\((G_{S'},T_{S'},\alpha_{S'})\) is an \(S'\)-elementary system for
every \(S'\to S\). If \((G,T,\alpha)\) is an \(S\)-elementary
system, then so is \((G,T,-\alpha)\).

If \(S\) is a scheme, \(G\) a reductive \(S\)-group, \(T\) a
maximal torus of \(G\), and \(\alpha\) a root of \(G\) with respect
to \(T\), then \((Z_\alpha,T,\alpha)\) is an \(S\)-elementary
system \textup{(\SGARefLink{sga3:xix:paragraph:3-9}{Exposé XIX, 3.9})}.

Let \((G,T,\alpha)\) be an \(S\)-elementary system. The invertible
module \(\mathfrak g^\alpha\) is canonically endowed with a
\(T\)-module structure. Hence the vector bundle
\(W(\mathfrak g^\alpha)\) is also a \(T\)-module. On the other hand,
the inner automorphisms by elements of \(T\) make \(G\) a group
equipped with a group of operators \(T\).
